Topologia:
Rodzinę $\mathcal{O}$ podzbiorów zbioru $X$ nazywamy **topologią w zbiorze $X$**, jeżeli:
1. $\emptyset \in \mathcal{O},\ X \in \mathcal{O}$,
2. $V_i \in \mathcal{O},\ i = 1, 2, \dots, n \implies V_1 \cap V_2 \cap \dots \cap V_n \in \mathcal{O}$,
3. $\{V_\alpha\}$ – dowolna rodzina zbiorów w $\mathcal{O} \implies \bigcup_\alpha V_\alpha \in \mathcal{O}$.

Przestrzeń $X$ z topologią $\mathcal{O}$, a formalnie parę $(X, \tau)$, nazywamy **przestrzenią topologiczną**, a elementy rodziny $\mathcal{O}$ – **zbiorami otwartymi w $X$**.
---
Algebra:
Rodzinę $\mathcal{M} \subset 2^X$ podzbiorów przestrzeni $X$ nazywamy **algebrą** (ciałem zbiorów), gdy:
1. $X \in \mathcal{M}$,
2. $A \in \mathcal{M} \Rightarrow A' \in \mathcal{M}$,
3. $A, B \in \mathcal{M} \Rightarrow A \cup B \in \mathcal{M}$.

Z powyższych warunków wynika zamkniętość na różnicę, skończone sumy i skończone przekroje.
---
Sigma-algebra:
Rodzinę $\mathcal{M} \subset 2^X$ podzbiorów przestrzeni $X$ nazywamy **$\sigma$-algebrą** (lub **$\sigma$-ciałem**), gdy:
1. $X \in \mathcal{M}$,
2. $A \in \mathcal{M} \Rightarrow A' \in \mathcal{M}$,
3. $A_n \in \mathcal{M},\ n = 1, 2, 3, \dots \Rightarrow \bigcup_{n=1}^\infty A_n \in \mathcal{M}$.

Z prawa De Morgana wynika także zamkniętość na przeliczalne przekroje.
---
Zbiór potęgowy:
Jeżeli $X$ jest przestrzenią, to możemy zdefiniować maksymalną rodzinę jego podzbiorów:
$$2^X := \{A : A \subset X\}.$$
$2^X$ nazywamy **zbiorem potęgowym**. Jeżeli $|X| = n$, to $2^X$ składa się z $2^n$ podzbiorów przestrzeni $X$.
---
Sigma-algebra generowana:
Niech $\mathcal{M} \subset 2^X$ będzie rodziną podzbiorów $X$. Rodzinę
$$\bigcap \{\mathcal{A} : \mathcal{M} \subset \mathcal{A},\ \mathcal{A} - \sigma\text{-algebra}\}$$
nazywamy **$\sigma$-algebrą generowaną przez rodzinę $\mathcal{M}$**. Będziemy ją oznaczać $\sigma(\mathcal{M})$ lub $\mathcal{M}^*$.

Własności:
1. jeśli $\mathcal{M}$ jest $\sigma$-algebrą, to $\mathcal{M}^* = \mathcal{M}$,
2. $(\mathcal{F}^*)^* = \mathcal{F}^*$,
3. jeśli $\mathcal{F}_1 \subset \mathcal{F}_2 \subset 2^X$, to $\mathcal{F}_1^* \subset \mathcal{F}_2^*$.
---
Przestrzeń mierzalna:
Jeżeli $\mathcal{M} \subset 2^X$ jest $\sigma$-algebrą, to parę $(X, \mathcal{M})$ nazywamy **przestrzenią mierzalną**, a elementy $\mathcal{M}$ – **zbiorami mierzalnymi**.
---
Zbiory borelowskie:
Załóżmy, że $(X, \mathcal{O}_X)$ jest przestrzenią topologiczną. Istnieje najmniejsza $\sigma$-algebra $\mathcal{B}_X$ w $X$ taka, że każdy otwarty podzbiór przestrzeni $X$ należy do $\mathcal{B}_X$.

Elementy rodziny $\mathcal{B}_X$ nazywamy **zbiorami borelowskimi w $X$**, tzn.
$$\mathcal{O}_X^* = \mathcal{B}_X.$$
---
Rozszerzony zbiór liczb rzeczywistych:
Rozszerzony zbiór liczb rzeczywistych ma postać:
$$\overline{\mathbb{R}} = [-\infty, +\infty] = \mathbb{R} \cup \{-\infty, +\infty\}.$$
Naturalne uporządkowanie zachowuje porządek na $\mathbb{R}$, przy czym $-\infty < x < +\infty$ dla każdego $x \in \mathbb{R}$.